% Chapter 7: Discussion

\section{Interpretation and Contribution}

This thesis adds a distributional layer to science mapping. A domain label tells us how a work is classified. A centroid gives the average position of a field. A projection gives a readable visual display. Structural morphology asks a different question: how the retained papers of a field are arranged around one another. The eight metrics describe this arrangement through global dispersion, local density, hubness, and spectral structure.

The empirical chapters support this view from three directions. The static comparison shows broad domain tendencies, but also large variation inside domains. The temporal analysis shows that field profiles change in measurable ways, without a single direction for all fields. The relational analysis shows that field pairs mostly move apart over time, although a substantial minority move closer. OpenAlex domains, nearest neighbors, bridges, and isolates help interpret these relations, but they do not replace the metric evidence.

The contribution is therefore not a new classification of fields. It is a reproducible way to compare field shape while keeping taxonomy, semantic position, temporal change, and pairwise relations separate. This separation matters because fields that are close on a map, fields that share a taxonomy, and fields with similar morphology are not necessarily the same cases. Centroid drift makes the distinction even clearer: a field can move in semantic location without changing its internal shape in the same way.

\section{Limits of the Evidence}

The evidence is conditional on the OpenAlex corpus, the primary-topic taxonomy, the balanced sampling design, SPECTER2, and the selected metrics. Coverage differs across fields, years, languages, document types, and metadata practices. Title--abstract embeddings compress scientific documents and miss many contextual features. The high-dimensional space also brings anisotropy, density variation, and hubness \parencite{aggarwal2001surprising,radovanovic2010hubness}. For this reason, the results should be read as controlled measurements, not as a full description of science.

The balanced design supports comparison, but it is not a census of world science. It also assigns each paper to one primary subfield, so papers with mixed disciplinary identities are simplified. The strongest claims are therefore comparative: one field is more dispersed than another under the fixed representation, one pair becomes closer in the selected profile distance, or temporal slope profiles cluster more clearly than static profiles.

The analysis does not explain why the observed changes occurred. Convergence and divergence are profile-distance terms, not claims about collaboration, citation flows, disciplinary merger, or epistemic separation. The metrics describe structural patterns in one controlled representation. They should not be read as value judgments about the quality or importance of fields.

\section{Implications and Future Validation}

Future work should first test model and corpus robustness. The same sampling plan, metric code, and aggregation rules can be rerun with alternative scientific embedding models, multilingual or full-text representations, and other bibliographic infrastructures. Changing one component at a time would show which findings are stable and which depend on representation or coverage.

External validation should compare morphology with evidence that was deliberately excluded from the metric construction: citation networks, journal structures, funding categories, author mobility, and expert assessments. Agreement would strengthen the interpretation of specific cases. Disagreement would also be useful, because it would show where geometric shape and other forms of scientific organization separate.

Temporal validation should focus on the strongest structural shifts, convergences, divergences, and centroid-drift cases. Five-year windows make the comparison stable, but they can smooth short shocks and database changes. Annual, rolling, and event-centered trajectories, combined with expert inspection of extreme cases, would test whether morphology highlights meaningful field transformations or mainly representation artifacts.
